\documentclass[12pt]{article}
\usepackage{amsmath}
\usepackage{geometry}
\geometry{margin=1in}

\title{Log Likelihood Explainer}
\author{Bryan Richlinski, with editing by Copilot}

\begin{document}
\maketitle

\section*{What Likelihood Is}

Likelihood is the probability of observing a particular outcome, assuming a fixed model of how the data is generated.  
%  
For example, if we assume a fair coin flip (i.e., each flip has a 0.5 probability of landing heads), then observing two heads in two flips has a likelihood of \( 0.5 \times 0.5 = 0.25 \).  
%  
However, if we instead assume that the probability of heads is 0.3, then the likelihood of observing two heads becomes \( 0.3 \times 0.3 = 0.09 \).  
%
For \emph{reasons} we say that a Likelyhood of no observation is always 1. 
%  
Importantly, this is not the same as the probability of the model given the observation.  
%  
That would require prior beliefs about the model itself, which we do not assume.  
%  
Confusing these two quantities is known as the prosecutor's fallacy—mistaking \( P(\text{model} \mid \text{data}) \) for \( P(\text{data} \mid \text{model}) \).  
%  


\section*{Log Likelihood}

Log likelihood is simply the likelihood function wrapped in the logarithm function.
%
The base of the logarithm is an open question to me right now, but it seems that that decision is often dependent upon the field of study. 
%
This transformation has two main advantages.  
%  
First, the logarithm is a strictly increasing function, so maximizing the log likelihood is equivalent to maximizing the likelihood itself.  
%  
Second, it simplifies arithmetic: ratios of probabilities become differences, \( \log(a/b) = \log a - \log b \), and joint probabilities become sums, \( \log(a \times b) = \log a + \log b \).  
%  
This makes optimization and numerical computation much easier.  
%  
As well the concept of no observation is much more intuitive, while the likelyhood of not observing something is set to 1, the log likelihood of 1 becomes 0. 
%  

\end{document}
